% !TeX root = ../main.tex
\section{Теоретическое обоснование значений}
Согласно заданию, необходимо исследовать зависимость мощности на выходе измерительного плеча направленного ответвителя (CPL OUT) от частоты при различных нагрузках на основном выходе (OUT): 
\begin{itemize}
    \item холостой ход (XX);
    \item согласованная нагрузка (50 Ом);
    \item ненагруженные аттенюаторы на 3~дБ, 6~дБ и 10~дБ.
\end{itemize}
Также требуется сделать вывод о влиянии номинала аттенюатора на уровень измеряемого сигнала. 

Дополнительно представлены результаты измерений коэффициента отражения для ненагруженных аттенюаторов номиналом 3~дБ, 6~дБ, 10~дБ и 20~дБ в диапазоне частот от 0 до 1000 МГц.

В данной схеме порт CPL OUT направленного ответвителя настроен на измерение мощности \textbf{отраженной волны}. 

Если к порту OUT подключается \textit{ненагруженный} аттенюатор с собственным затуханием $L_{\text{атт}}$ (выраженным в децибелах), то СВЧ-сигнал претерпевает следующие изменения:
\begin{enumerate}
    \item Падающая волна проходит через аттенюатор, её мощность уменьшается на величину $L_{\text{атт}}$.
    \item На конце аттенюатора находится холостой ход (отсутствие нагрузки), от которого сигнал отражается практически полностью. Коэффициент отражения по напряжению $\Gamma \approx 1$.
    \item Отраженная волна проходит через аттенюатор в обратном направлении, теряя еще $L_{\text{атт}}$ своей мощности.
\end{enumerate}

Таким образом, общие потери на отражение (Return Loss, $RL$) для ненагруженного аттенюатора равны его удвоенному номиналу:
$$ RL \approx 2 \cdot L_{\text{атт}} \text{ (дБ)} $$

Мощность сигнала $P_{\text{CPL}}$, попадающего в измерительное плечо, вычисляется как:
$$ P_{\text{CPL}} = P_{\text{пад}} - C - 2 \cdot L_{\text{атт}} $$
где $P_{\text{пад}}$ — мощность падающей волны, $C$ — переходное ослабление направленного ответвителя.

\begin{figure}[H]
    \centering
    \includegraphics[width=1\textwidth]{ris11.png}
    \caption{Зависимости подавления аттенюаторов от частоты}
    \label{fig:ris10}
\end{figure}

Анализируя представленный график зависимости коэффициента отражения от частоты, можно убедиться в справедливости выведенной формулы. Рассмотрим разницу между уровнями сигналов для разных аттенюаторов:

\begin{itemize}
    \item При переходе от аттенюатора 3~дБ (синий график) к 6~дБ (красный график) разница номиналов составляет $\Delta L = 3$~дБ. Ожидаемое падение уровня отраженного сигнала: $2 \cdot 3 = 6$~дБ. На графике (например, в районе 400 МГц) мы видим падение примерно с +3~дБ до -4~дБ, что дает разницу около 7~дБ (близко к теоретическим 6~дБ).
    \item При переходе от 6~дБ (красный график) к 10~дБ (желтый график) $\Delta L = 4$~дБ. Ожидаемое падение: $2 \cdot 4 = 8$~дБ. На графике уровень падает с -4~дБ до -13~дБ (разница $\sim 9$~дБ).
    \item При переходе от 10~дБ к 20~дБ (зеленый график) ожидается падение на $20$~дБ. Однако на графике падение составляет лишь порядка $10-15$~дБ в зависимости от частоты. Это физически объясняется \textbf{конечной направленностью ответвителя}: при слишком малых отраженных сигналах (аттенюатор 20~дБ снижает отражение на 40~дБ) доминирующим в плече CPL OUT становится паразитное прохождение падающей волны, и график упирается в шумовой порог/предел направленности измерительной установки.
\end{itemize}

